\icode{N-EXM-FI-L3-260625-001}
某同学利用如图(\(a\))所示的实验装置探究物体做直线运动时平均速度与时间的关系。让小车左端和纸带相连，右端用细绳跨过定滑轮和钩码相连。钩码下落，带动小车运动，打点计时器打出纸带。某次实验得到的纸带和相关数据如图(\(b\))所示。

\begin{figure}[htbp]
\begin{minipage}[c]{.49\textwidth}
    \begin{tikzpicture}[scale=1.2,font=\scriptsize]
        \draw[fill=orange!70] (0,-.1) rectangle(4,0) (.5,-.1) rectangle(.8,-1) (3.5,-.1) rectangle(3.2,-1);
        \draw[fill=black] (.4,0) rectangle(.5,.3);
        \draw (.4,.5)..controls(0,.5)and(-.2,.2)..(-.5,0);
        \filldraw[fill=gray,rotate around={-5:(3.9,0)}] (0.2,0) rectangle(3.9,.1);
        \filldraw[fill=white,rotate around={-5:(3.9,0)}](.3,.1) rectangle(.6,.25) (2,.2) rectangle(2.5,.35);
        \draw[purple,rotate around={-5:(3.9,0)}] (.6,.25)--(2,.25) (2.5,.25)--(4.1,.25);
        \filldraw[fill=white,rotate around={-5:(3.9,0)}] (2.1,.15) circle(.05) (2.4,.15) circle(.05);
        \filldraw[fill=white,rotate around={18:(3.9,0)}] (3.9,0.05) rectangle(4.15,.1);
        \draw (4.1,.13) circle(.1);
        \fill (4.1,.13) circle(.05);
        \draw[purple] (4.2,.13)--(4.2,-.5);
        \draw (4.1,-.5) rectangle(4.3,-.7);
        \node at(4.2,-.9) {钩码};
        \node at(2.3,.6)[rotate=-5]{小车};
        \node at(.5,.7)[rotate=-5]{打点计时器};
        \node at(-.3,.2)[left]{纸带};
        \node at(.8,.15){木板};
        \node at (2.1,-1) {图(\(a\))};
    \end{tikzpicture}
\end{minipage}
\begin{minipage}[c]{.49\textwidth}
    \begin{tikzpicture}[scale=1,font=\scriptsize]
        \draw (-.3,.5)--(6,.5) (-.3,-.5)--(6,-.5) (-.3,.5)--(-.4,.4)--(-.2,-.3)--(-.3,-.5) (6,.5)--(5.8,.4)--(6.2,-.3)--(6,-.5);
        \foreach \x in{0,.5,1.3,2.4,3.8,5.5}
          \draw (\x,.15)--(\x,.6);
        \foreach \x/\c in{0/A,.5/B,1.3/C,2.4/D,3.8/E,5.5/F}
          \fill (\x,0)node[below]{\(\c\)} circle(.05);
        \foreach \x/\y/\n in{0/.5/6.60,.5/1.3/8.00,1.3/2.4/9.40,2.4/3.8/10.90,3.8/5.5/12.40}
          \draw[|<->|] (\x,.8)--node[above]{\(\n\)}(\y,.8);
        \node at(5.3,-.7) {单位：\,\qty{}{cm}};
        \node at(3,-1){图(\(b\))};
    \end{tikzpicture}
\end{minipage}
\end{figure}

(1)已知打出图(\(b\))中相邻两个计数点的时间间隔均为\,\qty{0.1}{s}，以打出\(A\)点时小车位置为初始位置，将打出\(B\)、\(C\)、\(D\)、\(E\)、\(F\)各点时小车的位移\(\Delta x\)填到表中，小车发生对应位移所用时间和平均速度分别为\(\Delta t\)和\(\overline{v}\)。表中\(\Delta x_{AD} =\) \underline{\hbox to 4em{}}\,\qty{}{cm}，\(\overline{v}_{AD} =\)\underline{\hbox to 4em{}} \,\qty{}{cm/s}。

\begin{table}[ht]
    \centering
    \begin{tblr}{|c|c|c|c|c|c|}
    \hline
    位移区间&\(AB\)&\(AC\)&\(AD\)&\(AE\)&\(AF\)\\
    \hline
    \(\Delta x\)/\qty{}{cm}&6.60&14.60&\(\Delta x_{AD}\)&34.90&47.30\\
    \hline
    \(\overline{v}\)/(\qty{}{cm\cdot s^{-1}})&66.0&73.0&\(\overline{v}_{AD}\)&83.7&94.6\\
    \hline
    \end{tblr}
\end{table}

(2)根据表中数据得到小车平均速度\(\overline{v}\)随时间\(\Delta t\)的变化关系，如图(\(c\))所示。在答题卡上的图中补全实验点

\begin{figure}[htbp]
    \centering
    \begin{tikzpicture}[font=\scriptsize,>=latex]
    \draw[help lines,color=gray!80,very thin,step=.2] (0,0)grid(6,5);
    \draw[->,blue] (0,0)--(7,0)node[below]{\(\Delta t\)/\,\qty{}{s}};
    \draw[->,blue] (0,0)--(0,5.5)node[right]{\(\overline{v}\)/(\,\qty{}{cm\cdot s^{-1}})};
    \foreach \x/\y in {0/0,1/0.1,2/0.2,3/0.3,4/0.4,5/0.5,6/0.6}
        \node at(\x,0)[below,blue]{\(\y\)};
    \foreach \x/\y in {55/0,65/1,75/2,85/3,95/4,105/5}
        \node at(0,\y)[left,blue]{\(\x\)};
    \foreach \x/\y in{.8/1.1,2/1.8,4/3.23,5/3.9}
        \fill (\x,\y) circle(.05);
    \node at(3.2,-.8) {图(\(c\))};
    \end{tikzpicture}
\end{figure}
  
(3)从实验结果可知，小车运动的\(\overline{v} - \Delta t\)图线可视为一条直线，此直线用方程\(\overline{v} = k\cdot \Delta t + b\)表示，其中\(k =\)\underline{\hbox to 4em{}} \,\qty{}{cm/s^2}。(结果均保留3位有效数字)
  
(4)根据(3)中的直线方程可以判定小车做匀加速直线运动，得到打出\(A\)点时小车速度大小\(v_A =\)\underline{\hbox to 4em{}}，小车的加速度大小\(a =\)\underline{\hbox to 4em{}}。(结果用字母\(k\)、\(b\)表示)